\chapter{Tetrahedral residual-bundle origin gate}

Предыдущий cubic-root gate использовал transition
\(Z_{H,\nu}^3=I\), но ещё не объяснял, почему degree три является частью
project carrier, а не новым дискретным input. Здесь показывается, что
\(\mathbb Z_3\) уже содержится в symmetry-breaking chain выбранных
projectors. Одновременно возникает новая развилка: residual subgroup даёт
gauge bundle только при gauged или boundary-framed family symmetry.

\section{Канонический spin-three tetrahedral tensor}

Пусть \(B:\mathbb R^3\to u^\perp\subset\mathbb R^4\) --- standard triplet
isometry, а четыре projector axes равны
\begin{equation}
 n_a=B^T\operatorname{diag}\!\left[
 \frac{2}{\sqrt3}\left(P_a-\frac14I_4\right)
 \right],
 \qquad a=1,\ldots,4.
 \label{eq:project-tetrahedral-axes}
\end{equation}
Они удовлетворяют
\begin{equation}
 \|n_a\|=1,
 \qquad
 \sum_an_a=0,
 \qquad
 n_a\cdot n_b=-\frac13\quad(a\ne b).
\end{equation}
Следовательно, это вершины regular tetrahedron в family triplet.

Из них без нового continuous direction строится symmetric rank-three
tensor
\begin{equation}
 \mathcal T_{ijk}=\sum_{a=1}^4(n_a)_i(n_a)_j(n_a)_k.
 \label{eq:project-spin-three-tensor}
\end{equation}
Его trace равен нулю:
\begin{equation}
 \mathcal T_{iik}=\sum_a\|n_a\|^2(n_a)_k=0,
\end{equation}
поэтому \(\mathcal T\) лежит в real spin-three representation
\(\operatorname{Sym}^3_0(\mathbb R^3)\) dimension семь.

Even permutations четырёх axes реализуются proper rotations standard
triplet и сохраняют \(\mathcal T\). Тем самым
\begin{equation}
 \operatorname{Stab}_{SO(3)}(\mathcal T)=A_4.
 \label{eq:so3-to-a4-breaking}
\end{equation}
Project computation проверяет все двенадцать proper tetrahedral rotations;
максимальный tensor residual меньше \(1.8\times10^{-15}\).

Именно такая UV architecture теперь имеет прямое литературное
подтверждение. В \path{arXiv:2607.12366} построена renormalizable
\(SO(3)\) gauge theory со spin-three Higgs, breaking
\(SO(3)\to A_4\), и найдены finite-tension vortices с order-two и
order-three \(A_4\)-holonomies. Более ранний scalar-potential анализ
spin-three \(SO(3)\) field и tetrahedral vacuum приведён в
\path{arXiv:0910.4392}. Эти работы не выводят project finite algebra, но
показывают, что требуемый symmetry-breaking pattern допускает обычный
local gauge-Higgs UV completion.

\section{Projector оставляет ровно \texorpdfstring{\(\mathbb Z_3\)}{Z3}}

После выбора одного projector \(P_a\), или эквивалентно axis \(n_a\),
остаточный subgroup равен
\begin{equation}
 \operatorname{Stab}_{A_4}(n_a)
 =\{1,C_{a,+},C_{a,-}\}
 \simeq\mathbb Z_3.
 \label{eq:projector-z3-stabilizer}
\end{equation}
Это следует также из orbit--stabilizer identity:
\begin{equation}
 |A_4|=12,
 \qquad
 |\operatorname{Orb}(n_a)|=4,
 \qquad
 |\operatorname{Stab}(n_a)|=3.
\end{equation}
Два nontrivial elements являются ровно двумя inverse three-cycles,
выбранными winding \(\nu=\pm1\).

Machine audit дополнительно даёт
\begin{equation}
 [A_4,A_4]=V_4,
 \qquad
 A_4^{\rm ab}=A_4/V_4\simeq\mathbb Z_3.
\end{equation}
Это согласуется со старым binary-tetrahedral ledger:
\(2T^{\rm ab}\simeq\mathbb Z_3\). Direct replacement внешней geometry на
\(S^3/2T\) ранее закрыт, но internal gauge lift этому no-go не
противоречит.

Для каждого nontrivial residual rotation его lift в \(SU(2)\) имеет
порядок шесть:
\begin{equation}
 \widetilde C_{a,\nu}^{,3}=-1,
 \qquad
 \widetilde C_{a,\nu}^{,6}=1.
 \label{eq:binary-z6-lift}
\end{equation}
Следовательно, cubic-root transition предыдущей главы является local
frame description residual \(\mathbb Z_3\)-bundle, а не произвольно
добавленной третьей корневой ветвью.

\begin{theorem}[Residual origin of the cubic root]
Projector equation \(Q(H)=0\), дополненная orientation-preserving
tetrahedral carrier, имеет canonical residual subgroup
\(\operatorname{Stab}_{A_4}(P_a)\simeq\mathbb Z_3\). Его два nontrivial
holonomies совпадают с \(C_{a,\pm}\), а binary lift имеет order six.
Поэтому denominator \(1/3\) и transition \(Z^3=I\) имеют внутреннее
group-theoretic происхождение.
\end{theorem}

\section{Почему charge-two \texorpdfstring{\(B-L\)}{B-L} недостаточно}

Нельзя приписать этот результат одному abelian pairing Higgs. Condensate
поля charge \(n\) в обычной \(U(1)\)-теории оставляет residual
\(\mathbb Z_n\); continuum/TQFT derivation приведён, например, в
\path{arXiv:1307.4793}. Поэтому существующий charge-two
\(B-L\) pairing field оставляет только
\begin{equation}
 U(1)_{B-L}\longrightarrow\mathbb Z_2,
\end{equation}
а independent abelian origin \(\mathbb Z_3\) потребовал бы нового
charge-three condensate.

Project data указывают на более экономное разделение ролей:
\begin{equation}
 \boxed{
 \begin{aligned}
 \text{ordinary charge-two }B-L\text{ texture}
 &\quad\Rightarrow\quad
 \text{Majorana pairing и unit vortex},\\
 \text{tetrahedral family carrier}+P_a
 &\quad\Rightarrow\quad
 \text{residual }\mathbb Z_3\text{ frame holonomy}.
 \end{aligned}}
 \label{eq:two-role-root-architecture}
\end{equation}
Это устраняет необходимость charge-three scalar, но требует вывести оба
blocks из одного graded parent.

\section{Gauge/global fork}

Group-theoretic stabilizer ещё не равен physical gauge bundle.

Если \(A_4\) является только global family symmetry, четыре projector axes
могут быть физически различимы, но residual \(\mathbb Z_3\) не поставляет
connection и topological flux. Если \(A_4\) gauged как остаток
\(SO(3)_F\), order-three vortices и bundles существуют, но четыре axes
связаны gauge transformations. Тогда физически определена conjugacy class
holonomy, а конкретная family axis требует boundary framing или
дополнительного matter block, делающего relative orientation observable.

\begin{nogocriterion}[Tetrahedral gauge/global fork]
Residual identity~\eqref{eq:projector-z3-stabilizer} не является полным
family prediction, пока не определено, является ли tetrahedral symmetry
gauged, global или boundary-framed. Gauged branch должна объяснить
наблюдаемость relative axis; global branch должна независимо породить
connection.
\end{nogocriterion}

\section{Обновлённый decisive gate}

Discrete root degree теперь закрыт. Остаются три связанные задачи:
\begin{enumerate}
 \item вывести spin-three tensor \(\mathcal T\) или эквивалентный
 tetrahedral gauge reduction из finite graded algebra;
 \item получить boundary framing, связывающий residual holonomy с family
 matter без gauge-dependent axis assignment;
 \item из того же parent action вывести nonzero ordinary charge-two
 \(B-L\) defect condensate и его coupling к residual \(\mathbb Z_3\)
 frame.
\end{enumerate}

Таким образом, cubic-root bundle больше не является независимой загадкой.
Новый blocker --- не существование \(\mathbb Z_3\), а единое физическое
embedding tetrahedral gauge/frame sector и Majorana condensate.

Следующий gauge--family locking gate закрывает первую половину этого
blocker. Real bifundamental \(X\) ломает
\(SO(3)_{\rm gauge}\times SO(3)_{\rm family}^{\rm global}\) до exact
diagonal global group, поэтому gauge holonomy получает physical family
frame. Открытым остаётся уже только parent-derived pairing instability:
symmetry разрешает как condensed, так и uncondensed phase.

Машинный ledger сохранён в
\path{s2t_v4_family_defect_tetrahedral_residual_bundle_gate_results.json};
генератор ---
\path{s2t_v4_family_defect_tetrahedral_residual_bundle_gate.py}.